% !TeX program = xelatex
% Generated by DocVortex. Compile with: latexmk -xelatex <file>.tex
\documentclass[UTF8,fontset=fandol,11pt]{ctexart}
\usepackage[a4paper,margin=20mm]{geometry}
\usepackage{amsmath,amssymb}
\usepackage{graphicx}
\usepackage{longtable}
\usepackage{array}
\usepackage{booktabs}
\usepackage{multirow}
\usepackage{enumitem}
\usepackage[table]{xcolor}
\usepackage[normalem]{ulem}
\usepackage{xeCJKfntef}
\usepackage{fvextra}
\usepackage[unicode,hidelinks]{hyperref}
\definecolor{DocVortexGray}{HTML}{6B7280}
\definecolor{DocVortexSurface}{HTML}{F3F4F6}
\definecolor{DocVortexTableHeader}{HTML}{E5E7EB}
\setlength{\parindent}{2em}
\setlength{\parskip}{0.4em}
\setlist{leftmargin=*,nosep}
\setlength{\LTpre}{0.4em}
\setlength{\LTpost}{0.6em}
\newcommand{\DocVortexLineBreak}{\leavevmode\hfill\break}
\newcommand{\DocVortexPlaceholder}[1]{\fcolorbox{DocVortexGray}{DocVortexSurface}{\textit{#1}}}
\hypersetup{pdftitle={Throughput of Tiered Document Parsing}}
\begin{document}

\section*{Throughput of Tiered Document Parsing}

Ada Example, Bela Minta\par

\textbf{Abstract.} We compare four parsing tiers on a small synthetic corpus and report how layout analysis, optical character recognition and formula recognition change throughput and fidelity.\par

\subsection*{1. Introduction}

Let \(x_{i}\in\mathbb{R}^{d}\) denote the feature vector of page i, and let \(\alpha \in (0,1)\) be a mixing weight. The empirical loss is \(\mathcal{L}(\theta) = \frac{1}{N}\sum_{i = 1}^{N}\ell (f_{\theta}(x_{i}),y_{i})\) and the regulariser satisfies \(\| \theta \| _2\leq 1.^1\) The latency model that we fit is given by\par

\begin{equation*}
T (n) = a + b n + c n ^ {2}, \quad \int_ {0} ^ {\infty} e ^ {- x ^ {2}} d x = \frac {\sqrt {\pi}}{2}\tag{1}
\end{equation*}

where is the number of pages and , , are fitted constants. In Section 2 we report the measuredn a b c latencies, and we discuss the Hungarian-language test page in Section 3.\par

\subsection*{2. Method}

Each document is parsed once per tier. The median and the 95th percentile of the wall-clock time are recorded for every run, and the results are summarised in Table 1 on the next page.\par

{\footnotesize\color{DocVortexGray} 1 The norm constraint is the usual unit-ball regularisation; it is stated here only as an example footnote.\par}

{\small\itshape Table 1. Parsing latency by tier (milliseconds).\par}

\begin{longtable}{|>{\raggedright\arraybackslash}p{0.235000\linewidth}|>{\raggedright\arraybackslash}p{0.235000\linewidth}|>{\raggedright\arraybackslash}p{0.235000\linewidth}|>{\raggedright\arraybackslash}p{0.235000\linewidth}|}
\hline
Tier & Pages & Median & p95 \\
\hline
flash & 12 & 810 & 1 240 \\
\hline
basic & 12 & 1 126 & 2 980 \\
\hline
standard & 12 & 4 310 & 9 870 \\
\hline
advanced & 12 & 7 450 & 15 200 \\
\hline
\end{longtable}

\noindent\includegraphics[width=\linewidth,height=.8\textheight,keepaspectratio]{\detokenize{images/page_1_chart_3.jpg}}\par

{\small\itshape Figure 1. Relative throughput by tier.\par}

\subsection*{3. Hungarian test page}

Az árvíztűrő tükörfúrógép egy magyar pangram: minden ékezetes magyar betűt tartalmaz, így az ő és az ű betűt is. A dokumentum feldolgozása után a szövegnek pontosan meg kell maradnia. ÁRVÍZTŰRŐ TÜKÖRFÚRÓGÉP.\par

The quadratic term satisfies \(c > 0\) for every tier, and the fitted curve passes within \(\epsilon = 0.05\) of each measurement.\par

\end{document}
